% !TEX program = xelatex
\documentclass[11pt,a4paper]{article}
\usepackage{fontspec}
\setmainfont{Calibri}
\setmonofont[Scale=0.88]{Consolas}
\newfontfamily\symbolfont{Segoe UI Symbol}
\newfontfamily\cjkfont{Microsoft YaHei}[Scale=0.9]
\usepackage{polyglossia}\setdefaultlanguage{polish}
\usepackage[table]{xcolor}
\usepackage{booktabs,longtable,array,geometry,enumitem,graphicx,fancyhdr,caption}
\usepackage[most]{tcolorbox}
\PassOptionsToPackage{hyphens}{url}
\usepackage[hidelinks,unicode]{hyperref}
\emergencystretch=2.5em
\geometry{a4paper,margin=2.1cm,top=2.3cm,bottom=2.3cm}
\setlength{\parindent}{0pt}\setlength{\parskip}{5pt}
\definecolor{apteal}{HTML}{0E7C86}
\definecolor{apteallight}{HTML}{E6F4F5}
\definecolor{apgrey}{HTML}{F3F4F6}
\definecolor{apamber}{HTML}{B45309}
\definecolor{apamberlight}{HTML}{FEF3C7}
\definecolor{apred}{HTML}{B91C1C}
\definecolor{apredlight}{HTML}{FEE2E2}
\definecolor{apgreen}{HTML}{15803D}
\definecolor{apgreenlight}{HTML}{DCFCE7}
\renewcommand{\arraystretch}{1.25}
\captionsetup{font=small,labelfont=bf}
\pagestyle{fancy}\fancyhf{}
\fancyfoot[L]{\small AMPERE POINT — testy bez internetu · Q11 z modułem Shelly · v1 · 2026-09-28}
\fancyfoot[R]{\small\thepage}
\renewcommand{\headrulewidth}{0pt}
\newtcolorbox{uwaga}[1][Uwaga]{enhanced,breakable,colback=apamberlight,colframe=apamber,title={#1},fonttitle=\bfseries,boxrule=0.6pt,arc=2pt,left=6pt,right=6pt,top=4pt,bottom=4pt}
\newtcolorbox{info}[1][Jak to działa]{enhanced,breakable,colback=apteallight,colframe=apteal,title={#1},fonttitle=\bfseries,boxrule=0.6pt,arc=2pt,left=6pt,right=6pt,top=4pt,bottom=4pt}
\newtcolorbox{wazne}[1][Ważne]{enhanced,breakable,colback=apredlight,colframe=apred,title={#1},fonttitle=\bfseries,boxrule=0.6pt,arc=2pt,left=6pt,right=6pt,top=4pt,bottom=4pt}
\newtcolorbox{przyklad}[1][Przykład liczbowy]{enhanced,breakable,colback=apgreenlight,colframe=apgreen,title={#1},fonttitle=\bfseries,boxrule=0.6pt,arc=2pt,left=6pt,right=6pt,top=4pt,bottom=4pt}
\newtcolorbox{kodbox}{enhanced,breakable,colback=apgrey,colframe=apgrey,boxrule=0pt,arc=2pt,left=5pt,right=5pt,top=3pt,bottom=3pt,fontupper=\ttfamily\footnotesize}
\setlist{nosep,leftmargin=*}
\setcounter{secnumdepth}{2}
\begin{document}

{\LARGE\bfseries Testy bez internetu: Q11 z modułem Shelly\par}
\vspace{4pt}\textcolor{apteal}{\rule{\textwidth}{1.2pt}}
\emph{Poziom C: Wi-Fi i sieć lokalna są, internetu nie ma. Testy C1–C12 · AMPERE POINT · 28 września 2026 · oprac. Dawid Cekała}


\section*{W skrócie}

\begin{itemize}
\item \textbf{Poziomu C dotąd naprawdę nie sprawdziliśmy.} 25.09 internet był w tle, a jego brak symulowaliśmy tylko dla serwera czasu. Teraz moduł dostaje adres bez bramy do internetu: widzi laptopa, broker i serwer czasu, ale nie internet.
\item \textbf{Co już wiemy o czasie:} bez serwera czasu moduł wysyła sterownikowi zerową datę. Z lokalnym serwerem czasu poprawna godzina przychodzi około 18 s po starcie (25.09).
\item \textbf{Najważniejsze testy:}
\begin{itemize}
\item \textbf{C8, zanik zasilania:} czy moduł po starcie nie narzuci sterownikowi zapamiętanego limitu prądu;
\item \textbf{C9, restart samego modułu w trakcie ładowania:} tak będzie przy każdej aktualizacji modułu.
\end{itemize}
\item \textbf{Opcjonalnie:} C10 (serwer czasu podany przez router), C11 i C12 (porównanie z Tuya). Wszystkie trzy wymagają zmian w routerze, które ustawia administrator sieci.
\item Poziomy łączności, pojęcia i porównanie z Tuya są w dokumencie 00, w tabeli 3.2.
\end{itemize}

\section{Przygotowanie: moduł bez internetu}

\begin{longtable}{>{\raggedright\arraybackslash}p{3.00cm}>{\raggedright\arraybackslash}p{7.11cm}>{\raggedright\arraybackslash}p{5.17cm}}
\toprule
\rowcolor{apteallight}\textbf{Krok} & \textbf{Co robimy} & \textbf{Wynik oczekiwany} \\
\midrule\endhead
\bottomrule\endfoot
P1 & Sprawdzamy, że sieć modułu (punkt dostępowy) jest włączona & droga ratunkowa, gdyby moduł zniknął z sieci biura \\
P2 & Moduł dostaje stały adres 192.168.0.238, maskę sieci biura i \textbf{brak bramy}; serwer czasu 192.168.0.57 & moduł wraca pod tym samym adresem, monitor działa \\
P3 & Z modułu próba pobrania strony z internetu (polecenie pobrania adresu WWW przez moduł) & błąd: moduł nie ma internetu. Pobranie z laptopa na 192.168.0.57 działa \\
P4 & Na koniec wszystkich testów C: powrót na adres automatyczny i na publiczny serwer czasu & moduł znów ma internet, poziom B \\
\end{longtable}

\textbf{Narzędzie do napisania: serwer czasu z przesunięciem} (\texttt{narzedzia\textbackslash{}zegar\_\allowbreak{}przesuniety.\allowbreak{}py}, test C6).

\begin{itemize}
\item \textbf{Co robi:} podaje czas przesunięty tak, żeby za 2 minuty wybiła pełna godzina. Okno godzin ma tylko pełne godziny, więc bez tego test trwałby do dwóch godzin.
\item \textbf{Dlaczego zegar laptopa zostaje bez zmian:} moduł przekazuje sterownikowi tylko to, co dostał od serwera czasu, więc przesuwamy wyłącznie zegar ładowarki.
\item \textbf{Jak go uruchomić:} na czas testu zatrzymujemy kontener z serwerem czasu, bo zajmuje port 123. Po każdej zmianie czasu restartujemy sam moduł, żeby od razu pobrał nową godzinę.
\end{itemize}

\textbf{Laptop nie może usypiać podczas testów C.} 29.09 rano, po uśpieniu laptopa, serwer czasu w kontenerze przez około 48 min nie miał źródeł (log kontenera: „Can't synchronise: no selectable sources”). W tym czasie podawał czas z zegara maszyny wirtualnej, który po uśpieniu stoi w miejscu. DevKit przyjął godzinę spóźnioną o 44 minuty i przekazywał ją sterownikowi. Serwer czasu na laptopie jest dobry do testów, nie jako rozwiązanie u klienta.


\clearpage

\section{Testy}

\subsection{C1. Aplikacja Shelly w sieci lokalnej}

\begin{longtable}{>{\raggedright\arraybackslash}p{3.00cm}>{\raggedright\arraybackslash}p{4.98cm}>{\raggedright\arraybackslash}p{7.30cm}}
\toprule
\rowcolor{apteallight}\textbf{Krok} & \textbf{Co robimy} & \textbf{Wynik oczekiwany} \\
\midrule\endhead
\bottomrule\endfoot
C1a & Telefon w sieci biura, aplikacja Shelly; moduł bez internetu i bez chmury & czy aplikacja znajduje moduł w sieci lokalnej i steruje nim. Telefon ma w tym czasie internet przez sieć biura, więc sprawdzamy tylko drogę telefon {\symbolfont→} moduł \\
\end{longtable}

\subsection{C2. Sterownik sterowany z modułu}

\begin{longtable}{>{\raggedright\arraybackslash}p{3.00cm}>{\raggedright\arraybackslash}p{4.05cm}>{\raggedright\arraybackslash}p{8.23cm}}
\toprule
\rowcolor{apteallight}\textbf{Krok} & \textbf{Co robimy} & \textbf{Wynik oczekiwany} \\
\midrule\endhead
\bottomrule\endfoot
C2a & Symulator: wolne {\symbolfont→} podłączone {\symbolfont→} ładuje & sterownik przechodzi stany „wolna” {\symbolfont→} „auto podłączone” {\symbolfont→} „ładuje”, zamyka stycznik; moduł pokazuje to w polach \\
C2b & Limit prądu z modułu przy ładowaniu: 12 A & sterownik potwierdza w 1 s; wyświetlacz pokazuje 12 A \\
\end{longtable}

\subsection{C3. Godzina po starcie}

\begin{longtable}{>{\raggedright\arraybackslash}p{3.00cm}>{\raggedright\arraybackslash}p{4.90cm}>{\raggedright\arraybackslash}p{7.38cm}}
\toprule
\rowcolor{apteallight}\textbf{Krok} & \textbf{Co robimy} & \textbf{Wynik oczekiwany} \\
\midrule\endhead
\bottomrule\endfoot
C3a & Wyłączamy i włączamy ładowarkę; obserwujemy pierwsze ramki czasu & przez około 20 s zerowa data, potem poprawna; zapisujemy, co w tym czasie pokazuje wyświetlacz \\
\end{longtable}

\subsection{C4. Drogi lokalne bez internetu}

\textbf{Cel.} Potwierdzić, że wszystko, co sprawdziliśmy na poziomie B, działa też bez internetu. Ruch zostaje w sieci lokalnej, więc oczekujemy tego samego wyniku.


\begin{longtable}{>{\raggedright\arraybackslash}p{3.00cm}>{\raggedright\arraybackslash}p{5.57cm}>{\raggedright\arraybackslash}p{6.71cm}}
\toprule
\rowcolor{apteallight}\textbf{Krok} & \textbf{Co robimy} & \textbf{Wynik oczekiwany} \\
\midrule\endhead
\bottomrule\endfoot
C4a & Zmiana limitu przez HTTP i WebSocket & sterownik potwierdza \\
C4b & Zmiana limitu przez MQTT & sterownik potwierdza; broker działa w sieci lokalnej \\
C4c & Zmiana limitu przez most OCPP & sterownik potwierdza; most i serwer testowy OCPP w sieci lokalnej \\
C4d & Strona WWW z laptopa & działa jak w B2 \\
C4e & Home Assistant, jeśli B3 się udało & działa jak w B3 \\
\end{longtable}

\subsection{C5. Ikona sieci}

\begin{longtable}{>{\raggedright\arraybackslash}p{3.00cm}>{\raggedright\arraybackslash}p{4.91cm}>{\raggedright\arraybackslash}p{7.37cm}}
\toprule
\rowcolor{apteallight}\textbf{Krok} & \textbf{Co robimy} & \textbf{Wynik oczekiwany} \\
\midrule\endhead
\bottomrule\endfoot
C5a & Obserwacja stanu sieci wysyłanego do sterownika i ikony na wyświetlaczu & według opisu protokołu Tuya: 2 = brak routera, 3 = router bez chmury, 4 = połączony z chmurą. Zapisujemy, czy moduł bez internetu nadal wysyła „4”. Ciąg dalszy przy wyłączonym Wi-Fi w D8 \\
\end{longtable}

\subsection{C6. Okno godzin, tryb pracy i limit energii, z lokalnym czasem}

\textbf{Cel.} Te funkcje ma sam sterownik. Ustalić, czy je zgłasza i wykonuje, kiedy ma godzinę z lokalnego serwera. Przy okazji powstanie lista pól do dodania w produkcie Shelly.


\begin{longtable}{>{\raggedright\arraybackslash}p{3.00cm}>{\raggedright\arraybackslash}p{5.57cm}>{\raggedright\arraybackslash}p{6.71cm}}
\toprule
\rowcolor{apteallight}\textbf{Krok} & \textbf{Co robimy} & \textbf{Wynik oczekiwany} \\
\midrule\endhead
\bottomrule\endfoot
C6a & Serwer z przesunięciem: za 2 min wybije godzina G. Z menu: tryb „harmonogram”, okno od G do G+1. Symulator w stanie „ładuje” & sterownik zgłasza tryb „harmonogram” i okno (G, G+1); przed G stan „czeka”, o G ładowanie rusza \\
C6b & Przesunięcie na 2 min przed G+1 & o G+1 ładowanie staje \\
C6c & Z menu: tryb „energia”, limit 1 kWh & sterownik zgłasza tryb „energia” i, oby, limit energii; bez obciążenia licznik nie ruszy, więc sprawdzamy tylko zgłaszanie \\
\end{longtable}

\textbf{Ustalone 29.09 przy wgrywaniu v4.2 (poziom B):} zapis okna z modułu działa (2 bajty). Sterownik po przyjęciu okna \textbf{sam przechodzi w tryb „harmonogram”} i poza oknem przestaje ładować. Okna 0–0 nie przyjmuje. W C6 trzeba to uwzględnić: po teście ustawić tryb „natychmiast”.


\textbf{Zapis tych ustawień z modułu to osobny etap.} Moduł nie ma polecenia do zapisu dowolnego punktu danych; zapisuje tylko skrypt usługi, do pól zdefiniowanych w produkcie. Trzeba więc dodać pola w portalu i wgrać nową konfigurację, a to kasuje Wi-Fi modułu.


\subsection{C7. Harmonogram i webhook bez internetu}

\begin{longtable}{>{\raggedright\arraybackslash}p{3.00cm}>{\raggedright\arraybackslash}p{4.61cm}>{\raggedright\arraybackslash}p{7.67cm}}
\toprule
\rowcolor{apteallight}\textbf{Krok} & \textbf{Co robimy} & \textbf{Wynik oczekiwany} \\
\midrule\endhead
\bottomrule\endfoot
C7a & Harmonogram z B4 na poziomie C & zadania odpalają o czasie; godzina z lokalnego serwera \\
C7b & Webhook z B5 na poziomie C & wywołania dochodzą do laptopa jak w B5 \\
\end{longtable}

\subsection{C8. Zanik zasilania}

\textbf{Cel.} Rozstrzygnąć, co pamięta sterownik, co moduł i kto komu narzuca wartość po starcie.


\textbf{Co przewiduje kod naszego skryptu.} Po starcie skrypt czyta limit i włącznik ze sterownika i nadpisuje nimi pola w module. Do sterownika pisze tylko wtedy, gdy pole się zmieni na wartość inną niż ostatnio otrzymana ze sterownika. Normalnie wygrywa więc sterownik.


Słabe miejsce: zanim sterownik pierwszy raz zgłosi limit, skrypt nie zna jego wartości. Każda zmiana pola w tych kilku sekundach pójdzie do sterownika. Taką zmianą może być odtworzenie zapamiętanej wartości, harmonogram albo polecenie z aplikacji. Dlatego w C8d patrzymy na pierwsze sekundy logu.


\begin{longtable}{>{\raggedright\arraybackslash}p{3.00cm}>{\raggedright\arraybackslash}p{6.14cm}>{\raggedright\arraybackslash}p{6.14cm}}
\toprule
\rowcolor{apteallight}\textbf{Krok} & \textbf{Co robimy} & \textbf{Wynik oczekiwany / co zapisujemy} \\
\midrule\endhead
\bottomrule\endfoot
C8a & Stan wyjściowy z logu: limit 10 A ustawiony \textbf{z modułu}, ładowanie wyłączone, symulator „podłączone”. W drugiej rundzie dodatkowo tryb „harmonogram” z oknem, limit energii 1 kWh, wpis klucz-wartość i harmonogram z B4 & pełny zrzut wszystkich punktów danych i pól modułu \\
C8b & Wyłącznik ładowarki: wyłączyć na 30 s, włączyć & moduł i sterownik startują; monitor wraca \\
C8c & Zrzut po starcie porównany z C8a & które wartości sterownik zachował; oczekujemy limitu prądu, trybu, okna i limitu energii \\
C8d & Pierwsze sekundy logu po starcie & \textbf{czy moduł wysłał zapis limitu albo włącznika, zanim sterownik zgłosił swoje wartości} \\
C8e & Włącznik: przed wyłączeniem „włączone”, symulator „ładuje” & czy po powrocie zasilania ładowanie rusza samo, czy czeka na polecenie \\
C8f & Czas: po ilu sekundach od startu sterownik dostał poprawną godzinę & z logu \\
C8g & Moduł: Wi-Fi, stały adres, MQTT, serwer czasu, harmonogram, webhook, wpis klucz-wartość & wszystko na miejscu \\
C8h & Energia i czas sesji w polach modułu & energia sesji 0, bo bez obciążenia licznik nie rośnie. Czas sesji niewiarygodny: skrypt liczy go z zegara modułu, więc skacze przy każdej zmianie zegara. 29.09 przy korekcie zegara skoczył o 45 min (770 {\symbolfont→} 815), a po restarcie modułu na 1066 min, czyli mniej więcej tyle, ile moduł pracował przed restartem. Do poprawy w skrypcie: liczyć z czasu pracy modułu, nie z zegara \\
C8i & Powtórka C8a–C8d z limitem ustawionym \textbf{z menu ładowarki} & czy moduł po starcie nie nadpisze ustawienia z menu swoją zapamiętaną wartością \\
\end{longtable}

\textbf{Dlaczego to ważne.} Klient ustawia z menu 8 A, bo ma słabe zabezpieczenie. Moduł pamięta z aplikacji 16 A. Po zaniku zasilania moduł wstaje, wysyła 16 A i ładowarka rusza z 16 A. Taki błąd byłby groźny, dlatego C8d i C8i to najważniejsze punkty całego programu.


\textbf{Rekomendacja, jeśli C8c pokaże, że sterownik sam pamięta limit i włącznik.} Wyłączyć w portalu zapamiętywanie tych dwóch pól. Wtedy nic nie dają, bo sterownik i tak zgłasza swoje wartości po starcie, a są jedynym źródłem opisanego ryzyka.


\subsection{C9. Restart samego modułu w trakcie ładowania}

\textbf{Cel.} Sterownik pracuje dalej, moduł znika na kilkanaście sekund. Tak będzie przy każdej aktualizacji oprogramowania modułu.


\begin{longtable}{>{\raggedright\arraybackslash}p{3.00cm}>{\raggedright\arraybackslash}p{3.52cm}>{\raggedright\arraybackslash}p{8.76cm}}
\toprule
\rowcolor{apteallight}\textbf{Krok} & \textbf{Co robimy} & \textbf{Wynik oczekiwany / co zapisujemy} \\
\midrule\endhead
\bottomrule\endfoot
C9a & Symulator „ładuje”; restart modułu poleceniem & czy sterownik przerywa ładowanie, gdy moduł przestaje dawać znak życia \\
C9b & Po powrocie modułu & co moduł wysyła do sterownika; czy pola w module zgadzają się ze sterownikiem \\
\end{longtable}

\subsection{C10. Serwer czasu podany przez router, opcjonalnie}

\begin{longtable}{>{\raggedright\arraybackslash}p{3.00cm}>{\raggedright\arraybackslash}p{6.55cm}>{\raggedright\arraybackslash}p{5.73cm}}
\toprule
\rowcolor{apteallight}\textbf{Krok} & \textbf{Co robimy} & \textbf{Wynik oczekiwany} \\
\midrule\endhead
\bottomrule\endfoot
C10a & Administrator ustawia w routerze rozgłaszanie adresu serwera czasu (opcja 42 przy przydzielaniu adresów); w module pole serwera czasu puste; moduł na adresie automatycznym & czy moduł sam bierze ten serwer. Jeśli tak, w firmach godzina działałaby bez konfiguracji \\
\end{longtable}

\subsection{C11. Aplikacja Tuya bez internetu, opcjonalnie, porównanie}

\begin{longtable}{>{\raggedright\arraybackslash}p{3.00cm}>{\raggedright\arraybackslash}p{8.11cm}>{\raggedright\arraybackslash}p{4.17cm}}
\toprule
\rowcolor{apteallight}\textbf{Krok} & \textbf{Co robimy} & \textbf{Wynik oczekiwany} \\
\midrule\endhead
\bottomrule\endfoot
C11a & Q11 z modułem Tuya odcięta od internetu regułą w routerze; telefon w tej samej sieci; aplikacja Tuya & czy aplikacja steruje ładowarką lokalnie \\
\end{longtable}

\subsection{C12. Harmonogram Tuya bez internetu, opcjonalnie, porównanie}

\begin{longtable}{>{\raggedright\arraybackslash}p{3.00cm}>{\raggedright\arraybackslash}p{7.27cm}>{\raggedright\arraybackslash}p{5.01cm}}
\toprule
\rowcolor{apteallight}\textbf{Krok} & \textbf{Co robimy} & \textbf{Wynik oczekiwany} \\
\midrule\endhead
\bottomrule\endfoot
C12a & Wpis harmonogramu w aplikacji Tuya za 10 min; potem ładowarka odcięta od internetu regułą w routerze & jeśli wpis się wykona, wykonuje go moduł Tuya, a nie chmura \\
\end{longtable}

\section{Kolejność i czas}

\begin{longtable}{>{\raggedright\arraybackslash}p{3.00cm}>{\raggedright\arraybackslash}p{7.32cm}>{\raggedright\arraybackslash}p{4.96cm}}
\toprule
\rowcolor{apteallight}\textbf{Kolejność} & \textbf{Testy} & \textbf{Czas} \\
\midrule\endhead
\bottomrule\endfoot
1 & P1–P3, C1–C5 & 1 h 25 min \\
2 & C8 pierwsza runda, C9 & 1 h 5 min \\
3 & C6 z narzędziem, C7 & 1 h 20 min \\
4 & C8 druga runda & 45 min \\
5 & P4 powrót na poziom B & 10 min \\
— & C10–C12, opcjonalnie & gdy administrator ustawi router \\
\end{longtable}

Razem około 4 h 45 min. Miejsce tych testów w kolejności całego programu: dokument 00, rozdział 5.


\section{Wyniki}

\begin{longtable}{>{\raggedright\arraybackslash}p{3.11cm}>{\raggedright\arraybackslash}p{2.00cm}>{\raggedright\arraybackslash}p{8.57cm}>{\raggedright\arraybackslash}p{1.16cm}}
\toprule
\rowcolor{apteallight}\textbf{Test} & \textbf{Wynik} & \textbf{Co zaobserwowano} & \textbf{Data} \\
\midrule\endhead
\bottomrule\endfoot
P3 moduł bez internetu &  &  &  \\
Przygotowanie P1–P3 & \leavevmode{\symbolfont\color{apgreen}\textbf{✔}} & 30.09 13:56: moduł na stałym adresie 192.168.0.238 bez bramy i bez serwera nazw. Z modułu: nazwa strony w internecie „Host not found”, adres IP w internecie „Timed out”; laptop (192.168.0.57:8099) odpowiada. Punktu dostępowego nie włączaliśmy (jest otwarty, bez hasła); droga ratunkowa: ten sam adres i podsieć & 30.09 13:57 \\
C1 aplikacja w sieci lokalnej & \leavevmode{\symbolfont\color{apred}\textbf{✘}} z B1 & Wynika z B1: aplikacja steruje tylko przez chmurę & 30.09 \\
C2 sterownik z modułu & \leavevmode{\symbolfont\color{apgreen}\textbf{✔}} z uwagą & 30.09 14:57, bez internetu. C2a: tester B {\symbolfont→} C; moduł pokazał „ładuje” i sygnał 6 V, webhooki włącznik „wł.”, „ładuje”, start ładowania. Stanu „podłączone” (B) sterownik nie zgłosił, tak jak w B5b. C2b: limit 12 A z modułu przy ładowaniu — sterownik potwierdził w 1 s (najpierw stare 8, potem 12), wyświetlacz pokazał 12 A & 30.09 14:58 \\
C3 godzina po starcie & \leavevmode{\symbolfont\color{apamber}\textbf{◐}} & 30.09 14:06, zanik zasilania całej ładowarki, moduł bez internetu: sterownik dostał poprawną godzinę 23 s po włączeniu (z serwera czasu na laptopie). Co pokazuje wyświetlacz przez te sekundy: do sprawdzenia. Uwaga S-45: bez serwera czasu moduł ma do tego czasu starą, zapamiętaną godzinę & 30.09 14:07 \\
C4 drogi lokalne & \leavevmode{\symbolfont\color{apgreen}\textbf{✔}} & 30.09 14:08, bez internetu: zapis limitu przez HTTP (210 ms), WebSocket (150 ms) i MQTT przez broker na laptopie; sterownik potwierdził 9, 10 i 8 A w 1–2 s. Strona WWW odpowiada, HA bez błędów. C4c (most OCPP): \emph{korekta 01.10} — most jest (\texttt{DevKit\textbackslash{}ocpp\textbackslash{}most\_\allowbreak{}q11.\allowbreak{}py}, z 25.09), patrz C4c niżej & 30.09 14:09 \\
C4c most OCPP bez internetu & \leavevmode{\symbolfont\color{apgreen}\textbf{✔}} & 01.10 11:34–11:36, moduł bez internetu (zapytania do internetu kończą się po czasie). Most i testowy serwer OCPP 1.6J na laptopie (\texttt{ws:\allowbreak{}/\allowbreak{}/\allowbreak{}localhost:\allowbreak{}9000}). BootNotification: Accepted, StatusNotification: Available, heartbeat co 60 s. Z serwera: limit 10 A (SetChargingProfile: Accepted, sterownik potwierdził 10 A), TriggerMessage status i pomiary (Accepted, pomiary 0, bo tester odłączony), GetConfiguration, limit 12 A (potwierdzony). 25.09 ten sam most działał z modułem mającym internet (tylko chmura wyłączona). Serwera OCPP w internecie nie testowaliśmy. Pełna transakcja \leavevmode{\symbolfont\color{apgreen}\textbf{✔}} (11:36–11:39, tester A {\symbolfont→} C {\symbolfont→} A): StatusNotification „SuspendedEVSE”, potem „Charging” i StartTransaction nr 1 (licznik 0 Wh), MeterValues co 30 s (6 razy, L1 225 V, 0 A), po odłączeniu „Available” i StopTransaction z powodem „EVDisconnected” & 01.10 11:40 \\
C5 ikona sieci & \leavevmode{\symbolfont\color{apamber}\textbf{◐}} obserwacja & 30.09 14:5x, bez internetu: moduł co 30 s wysyła sterownikowi stan sieci 4 („połączony z chmurą”), bez zmian. Ikona Wi-Fi na wyświetlaczu co kilkadziesiąt sekund znika i wraca (zdjęcia Dawida). [Z] Sterownik gasi ikonę między meldunkami. Do sprawdzenia: to samo z internetem i po zmianie zgłaszanego stanu na 3 (rejestr S-39) & 30.09 14:57 \\
C6 okno godzin z lokalnym czasem & \leavevmode{\symbolfont\color{apgreen}\textbf{✔}} & 30.09 15:24–15:31, bez internetu, tester w C. Serwer czasu z przesunięciem (\texttt{narzedzia\textbackslash{}zegar\_\allowbreak{}przesuniety.\allowbreak{}py}, w kontenerze zamiast \texttt{ntp}; zmiana czasu = restart modułu). Okno 22–23 ustawione z modułu (w menu nie ma okna, tylko opóźnienie): sterownik przeszedł w „harmonogram” i przerwał ładowanie („czeka”). C6a: o 22:00:03 zegara modułu „ładuje”. C6b: po przestawieniu na 22:58 ładowanie trwa, o 23:00:03 „czeka”. \textbf{Sterownik wykonuje okno sam, według godziny, którą dostaje od modułu} (odpowiada mu na pytanie o czas co około 18 s). C6c pominięty: sprzężenie limitu energii z trybem znane z B (S-21). Po teście: zwykły serwer czasu, tryb „od razu”, ładuje & 30.09 15:31 \\
C7 harmonogram i webhook & \leavevmode{\symbolfont\color{apgreen}\textbf{✔}} & 30.09 14:11, bez internetu: zadanie harmonogramu odpaliło co do sekundy (godzina z serwera na laptopie), sterownik potwierdził 11 A; webhooki na zmianę limitu doszły do laptopa (14:08). Zadanie usunięte, limit z powrotem 8 A. Uwaga S-45: po zaniku zasilania bez serwera czasu harmonogram działa według starej godziny & 30.09 14:11 \\
C8 zanik zasilania & \leavevmode{\symbolfont\color{apamber}\textbf{◐}} runda 1 & 30.09 14:06, bez internetu. Limit ustawiony z menu (8 A): sterownik go zachował, moduł też; przy starcie moduł nie wysłał sterownikowi żadnego zapisu (log UDP od 21 s). Łącze gotowe po 30 s, skrypt bez komunikatów, godzina 23 s po starcie. Test pamięci (dokument 04): moduł odtwarza wszystko, co zapisuje; sterownik traci okno, tryb harmonogramu i opóźnienie z menu (S-47). C8e \leavevmode{\symbolfont\color{apgreen}\textbf{✔}} (15:19, tester w C): po zaniku zasilania ładowanie ruszyło samo — sterownik około 30 s po włączeniu zgłosił „ładuje”, limit 12 A, tryb „od razu”; moduł nic nie zapisał, godzinę przekazał po 24 s. Razem z S-47: po mrugnięciu prądu ładowarka ładuje od razu, nawet gdy było ustawione opóźnienie. Zostaje druga runda & 30.09 15:20 \\
C9 restart modułu w trakcie ładowania & \leavevmode{\symbolfont\color{apgreen}\textbf{✔}} z uwagą & 30.09 15:00–15:07, bez internetu, tester w C: 5 restartów modułu (każdy około 30 s bez modułu, trzy jeden po drugim). Sterownik ładował bez przerwy, na ekranie nic się nie działo (Dawid). Po powrocie moduł nie wysłał zapisu, pola zgodne ze sterownikiem (ładuje, 6 V, 12 A). Uwaga: około 40 s po restarcie przychodzi fałszywy webhook „ładuje” i „start ładowania” (rejestr S-50) & 30.09 15:07 \\
C10 serwer czasu z routera &  &  &  \\
C11 aplikacja Tuya &  &  &  \\
C12 harmonogram Tuya &  &  &  \\
\end{longtable}

\end{document}